\documentclass{beamer}
\usetheme{Madrid}
\usecolortheme{default}

\usepackage[utf8]{inputenc}
\usepackage{graphicx}
\usepackage{xcolor}

% Custom colors for pillars
\definecolor{abstraction}{RGB}{0, 102, 204}    % Blue
\definecolor{encapsulation}{RGB}{0, 153, 0}    % Green
\definecolor{inheritance}{RGB}{204, 102, 0}    % Orange
\definecolor{polymorphism}{RGB}{153, 0, 153}   % Purple

% Tight spacing
\setlength{\parskip}{0pt}
\setlength{\itemsep}{0pt}
\setlength{\parsep}{0pt}
\setlength{\topsep}{0pt}

\title{Object Oriented Programming in Python}
\subtitle{A Demonstration of the Four Pillars}
\author{Holyweek Asea}
\institute{Student Number: 2500721983 \\[0.1cm] Makerere University}
\date{September 1, 2026}

\begin{document}
	
	% ======================== SLIDE 1: COVER ========================
	\begin{frame}
		\titlepage
	\end{frame}
	
	% ======================== SLIDE 2: INTRODUCTION ========================
	\begin{frame}{Introduction to OOP \& Why Python?}
		
		\textbf{What is Object Oriented Programming?}
		\begin{itemize}
			\item OOP models real-world entities as \textbf{objects}.
			\item Each object has \textbf{attributes} (data) and \textbf{methods} (behavior).
			\item Organizes code to be reusable, maintainable, and scalable.
		\end{itemize}
		
		\vspace{0.1cm}
		
		\textbf{Why Python for OOP?}
		\begin{itemize}
			\item \textbf{Readable and Beginner Friendly}: Clear, simple syntax.
			\item \textbf{Supports All OOP Principles}: Classes, inheritance, polymorphism, encapsulation, and abstraction are built-in.
			\item \textbf{Widely Used in Statistics}: Python is the primary language for data science and statistical computing.
		\end{itemize}
		
	\end{frame}
	
	% ======================== SLIDE 3: CODE - CLASS DEFINITIONS ========================
	\begin{frame}[fragile]{School Grading System: Class Definitions}
		
		{\fontsize{5.5}{6}\selectfont
			\begin{verbatim}
				from abc import ABC, abstractmethod
				
				class Person(ABC):                        # 1. ABSTRACTION
				def __init__(self, name, id_no):      	  #    Abstract base class hides implementation details.
				self.name = name                  		      
				self._id_no = id_no                 
				
				@abstractmethod                           #    Forces children to implement this method.
				def display_details(self): pass          
				
				class Student(Person):                    # 2. INHERITANCE
				def __init__(self, name, id_no, marks):   #    Reuse parent constructor
				super().__init__(name, id_no)     
				self.__marks = marks                      # 3. ENCAPSULATION
				
				def get_marks(self):                      #    Private attribute accessed only via getter method.
				return self.__marks                   
														   
				
				def display_details(self):                # 4. POLYMORPHISM
				print(f"Student: {self.name}, Marks: {self.__marks}")
				#    Same method name,
				def calculate_grade(self):                #    different behavior for Student vs Teacher.
				marks = self.get_marks()            
				if marks >= 90: return "A"
				elif marks >= 70: return "B"
				elif marks >= 50: return "C"
				else: return "F"
				
				class Teacher(Person):                    # 2. INHERITANCE
				def __init__(self, name, id_no, subject):
				super().__init__(name, id_no)             #    Reuse parent constructor
				self.subject = subject
				
				def display_details(self):                # 4. POLYMORPHISM
				print(f"Teacher: {self.name}, Subject: {self.subject}") #    Same method name, different output.
				
			\end{verbatim}
		}
		
	\end{frame}
	
	% ======================== SLIDE 4: CODE - DEMONSTRATION & OUTPUT ========================
	\begin{frame}[fragile]{School Grading System: Demonstration \& Output}
		
		{\fontsize{6}{6.5}\selectfont
			\begin{verbatim}
				# Create instances
				student1 = Student("Holyweek Asea, Marks", 501, 88)
				teacher1 = Teacher("Faith Kuli", 101, "Computer Science")
				
				# 4. POLYMORPHISM in action
				people = [student1, teacher1]
				for p in people:
				p.display_details()      # Same method call, different output
				
				# 3. ENCAPSULATION in action
				print(f"Marks: {student1.get_marks()}")   # Access via getter
				
				# 1. ABSTRACTION in action
				print(f"Grade: {student1.calculate_grade()}")  # User doesn't need
				# to know HOW it works.
			\end{verbatim}
		}
		
		\vspace{0.1cm}
		
		\textbf{Output:}
		{\fontsize{6}{6.5}\selectfont
			\begin{verbatim}
				Student: Holyweek Asea, Marks: 88
				Teacher: Faith Kuli, Subject: Computer Science
				Marks: 88
				Grade: B
			\end{verbatim}
		}
		
	\end{frame}
	
	% ======================== SLIDE 5: SUMMARY + REFERENCES ========================
	\begin{frame}{The Four Pillars Demonstrated}
		
		\begin{center}
			\begin{tabular}{|p{0.22\textwidth}|p{0.65\textwidth}|}
				\hline
				\textbf{Pillar} & \textbf{Implementation in Code} \\
				\hline
				\hline
				{\color{abstraction}\textbf{Abstraction}} & \texttt{Person(ABC)} and \texttt{@abstractmethod} hide implementation details. \\
				\hline
				{\color{encapsulation}\textbf{Encapsulation}} & \texttt{self.\_\_marks} is \textbf{private}; accessed only via \texttt{get\_marks()}. \\
				\hline
				{\color{inheritance}\textbf{Inheritance}} & \texttt{Student(Person)} and \texttt{Teacher(Person)} inherit from \texttt{Person}. \\
				\hline
				{\color{polymorphism}\textbf{Polymorphism}} & \texttt{for p in people: p.display\_details()} behaves differently for each object. \\
				\hline
			\end{tabular}
		\end{center}
		
		\vspace{0.1cm}
		
		\noindent \textbf{References:}
		\vspace{0.05cm}
		
		\begin{thebibliography}{9}
			\bibitem{kaur2026} 
			Kaur, S. (2026). \textit{Mastering Object-Oriented Programming with Python}. Unpublished Textbook. Makerere University.
			
			\bibitem{python} 
			Python Software Foundation. (2026). \textit{Python 3 Documentation}. 
			Retrieved from https://docs.python.org/3/
		\end{thebibliography}
		
	\end{frame}
	
\end{document}